% generated by roam.py handbook export -- не редактировать, править заметки
% 20260922164226-l_гладкие_функции.org (body 20f2e75c3d8b)

\begin{definition}[{L-гладкие функции}]\label[definition]{note:7E287F4B-AC3E-4CF5-9D32-3BC2BC52D01A}
  Пусть $f: E \rightarrow \R$, $f \in C^n$ ($n$ раз \hyperref[note:957FA09D-8883-40A5-B4C7-237D9F815A42]{дифференцируема}), $L > 0$, $k \le n$.
  Тогда мы говорим, что  $f$  имеет $L$-гладкие $k$-е производные, если
  \[
  \forall x,y \in E\ \| \nabla^k f(x) - \nabla^k f(y) \|_2 \le L \|x - y\|_2
  \]
  Альтернативное обозначение: $f \in C^{n,k}_L$
\end{definition}

\begin{remark}
  В частности при $k=1$, $f$~--- \hyperref[note:7DD1C7AB-1A49-4675-9F21-5A16C717F670]{функция с липшицевым градиентом} ($f \in C^{n,1}_L$)
\end{remark}

\begin{remark}
  Норма берется по типу $\nabla^k f$: для $k = 1$ (градиент~--- вектор)~--- евклидова, для $k = 2$ (гессиан~--- матрица)~--- спектральная. Для $f \in C^2$ условие при $k = 1$ равносильно $\| \nabla^2 f(x) \| \le L$: см. \hyperref[note:4B6D53BA-E599-4DEA-A557-550BD32042BA]{связь спектральной нормы гессиана и липшицевости градиента} и \hyperref[note:DAA0DAE4-4B48-4A2F-9459-54790B6B129F]{оценку нормы гессиана сверху}.
\end{remark}
